Trong bối cảnh thương mại điện tử phát triển mạnh mẽ tại Việt Nam, nhu cầu mua
sắm trực tuyến ngày càng tăng cao, đặc biệt trong lĩnh vực thời trang và phụ kiện.
Tuy nhiên, nhiều doanh nghiệp vừa và nhỏ vẫn chưa có nền tảng bán hàng trực
tuyến riêng, phụ thuộc hoàn toàn vào các sàn thương mại điện tử trung gian như
Shopee hay Lazada, dẫn đến việc bị động trong quản lý đơn hàng, thông tin khách
hàng, và chi phí hoa hồng cao. Các giải pháp xây dựng website thương mại điện
tử hiện có hoặc quá phức tạp, tốn chi phí vận hành lớn, hoặc thiếu các nghiệp vụ
đặc thù như quản lý hoàn trả hàng, cấu hình phí vận chuyển theo địa bàn, hay phân
quyền linh hoạt cho nhân viên.
Trước thực tế đó, đồ án lựa chọn hướng tiếp cận xây dựng một ứng dụng web
thương mại điện tử hoàn chỉnh theo kiến trúc client-server, sử dụng bộ công nghệ
MERN Stack (MongoDB, Express.js, React.js, Node.js) — một lựa chọn phổ biến,
có hệ sinh thái mạnh, phù hợp với quy mô dự án vừa và khả năng mở rộng sau này.
Giải pháp được xây dựng là website bán hàng, bao gồm hai phần chính: giao
diện người dùng cho phép khách hàng duyệt sản phẩm, quản lý giỏ hàng, thanh
toán qua VNPay hoặc COD, xem lịch sử đơn hàng và yêu cầu hoàn trả; phần quản
trị hỗ trợ quản lý sản phẩm, danh mục, mã giảm giá, phân quyền nhân viên theo
vai trò, quản lý người dùng, đơn hàng, đổi trả và nhật ký hành động. Hệ thống tích
hợp thêm Socket.IO cho chatbot hỗ trợ khách hàng theo thời gian thực và gửi email
xác nhận qua Gmail OAuth2.
Đóng góp chính của đồ án là xây dựng thành công một hệ thống thương mại
điện tử đầy đủ nghiệp vụ thực tế, với quy trình hoàn trả hàng có vòng đời trạng thái
rõ ràng, hệ thống phân quyền chi tiết theo vai trò, và cấu hình phí vận chuyển linh
hoạt — đáp ứng nhu cầu vận hành thực tiễn của một cửa hàng thời trang trực tuyến
quy mô vừa.